56
10 *Differential Calculus from a More General Point of View
Proposition 1 For a multilinear transformation $A : X_1 \times \cdots \times X_n \to Y$ mapping a
product of normed spaces $X_1, \ldots, X_n$ into a normed space $Y$ the following condi-
tions are equivalent:
a) $A$ has a finite norm,
b) $A$ is a bounded transformation,
c) $A$ is a continuous transformation,
d) $A$ is continuous at the point $(0, \ldots, 0) \in X_1 \times \cdots \times X_n.$
Proof We prove a closed chain of implications $a) \Rightarrow b) \Rightarrow c) \Rightarrow d) \Rightarrow a).$
It is obvious from relation (10.27) that $a) \Rightarrow b).$
Let us verify that $b) \Rightarrow c),$ that is, that (10.29) implies that the operator $A$ is
continuous. Indeed, taking account of the multilinearity of $A,$ we can write that
$$A(x_1+ h_1, x_2 + h_2, \ldots, x_n + h_n) – A(x_1, x_2, \ldots, x_n) =$$
$$= A(h_1, x_2, \ldots, x_n) + \ldots + A(x_1, x_2, \ldots, x_{n-1}, h_n) =$$
$$+ A(h_1,h_2, x_3, \ldots, x_n) + \ldots + A(x_1,\ldots, x_{n-2}, h_{n-1}, h_n) +$$
$$+ \cdots +A(h_1,\ldots, h_n).$$
From (10.29) we now obtain the estimate
$$|A(x_1+ h_1, x_2 + h_2, \ldots, x_n + h_n) - A(x_1, x_2,\ldots,x_n) | \le$$
$$\le M(|h_1|\cdot|x_2|\cdots|x_n| + \cdots +|x_1|\cdot|x_2|\cdots |x_{n-1}|\cdot|h_n|+$$
$$\left. + \cdots + |h_1|\cdots |h_n| \right),$$
from which it follows that $A$ is continuous at each point $(x_1, \ldots, x_n) \in X_1 \times \cdots \times$
$X_n.$
In particular, if $(x_1, \ldots, x_n) = (0, \ldots, 0)$ we obtain $d)$ from $c).$
It remains to be shown that $d) \Rightarrow a).$
Given $\varepsilon > 0$ we find $\delta = \delta(\varepsilon) > 0$ such that $|A(x_1,\ldots,x_n)| < \varepsilon$ when $\max\{|x_1|,$
$\ldots, |x_n|\} < \delta.$ Then for any set $e_1, \ldots, e_n$ of unit vectors we obtain
$$\frac{1}{\delta^n} |A(e_1,\ldots, e_n)| = \left|A\left(\frac{\delta e_1}{\delta}, \ldots, \frac{\delta e_n}{\delta}\right)\right| < \frac{\varepsilon}{\delta^n},$$
that is, $||A|| < \infty.$
We have seen above (Example 1) that not every linear transformation has a finite
norm, that is, a linear transformation is not always continuous. We have also pointed
out that continuity can fail for a linear transformation only when the transformation
is defined on an infinite-dimensional space.